\documentclass[11pt,a4paper]{article}

% Encoding and fonts
\usepackage[utf8]{inputenc}
\usepackage[T1]{fontenc}
\usepackage{times}
\usepackage{latexsym}
\usepackage{inconsolata}
\usepackage{microtype}

% Graphics
\usepackage{graphicx}
\usepackage{float}

% Mathematics
\usepackage{amsmath,amssymb}

% Hyperlinks and cross-references
\usepackage{hyperref}
\hypersetup{colorlinks=true,linkcolor=blue,citecolor=blue}
\usepackage{cite}

% For proper table formatting
\usepackage{array,booktabs}

\title{Deep Residual Learning for Image Recognition: A Comprehensive Analysis}
\author{Anonymous}
\date{}

\begin{document}

\maketitle

\begin{abstract}
Deep convolutional neural networks have driven remarkable progress in image recognition, with the prevailing belief that increased depth leads to better accuracy. However, a critical obstacle known as the degradation problem emerged: as networks become deeper, their training error paradoxically increases, indicating optimization difficulties rather than overfitting. This paper provides a thorough analysis of the seminal ResNet architecture (He et al., 2016), which introduced residual learning to overcome this barrier. By reformulating the desired mapping as $H(x) = F(x) + x$, where $F(x)$ represents a residual function, the network learns perturbations to an identity mapping, substantially simplifying the optimization of ultra-deep models. We examine the motivation, architectural design, experimental results, and theoretical insights of ResNet, covering both the original paper and subsequent extensions. Through detailed explanations of shortcut connections, bottleneck blocks, and gradient flow, we demonstrate why residual learning enables networks as deep as 152 layers to achieve state-of-the-art performance on ImageNet with a top-5 error of 3.57\%. Furthermore, we discuss the broader impact of ResNet on subsequent architectures such as ResNeXt, DenseNet, and Wide ResNet, and its cross-domain adoption in natural language processing, speech, and reinforcement learning. This survey serves as a comprehensive resource for understanding the principles behind ResNet and its enduring influence on deep learning.
\end{abstract}

\section{Introduction}
\label{sec:intro}

Deep convolutional neural networks (CNNs) have been the cornerstone of computer vision since AlexNet (Krizhevsky et al., 2012) demonstrated their superiority on ImageNet. In the years that followed, researchers consistently observed that increasing network depth improved performance: VGG (Simonyan \& Zisserman, 2015) pushed to 19 layers, and GoogLeNet (Szegedy et al., 2015) introduced inception modules with 22 layers. The underlying intuition is straightforward — deeper networks can represent more abstract features and model more complex functions. However, a surprising phenomenon surfaced as depth was further increased: when plain networks (without residual connections) grew beyond a certain point, their training error began to rise, even though model capacity was larger. This contradicted the expectation that deeper models should at least achieve the same training loss as shallower ones (since a deeper net can simply learn identity mappings for additional layers).

The degradation problem is distinct from overfitting or vanishing gradients. While vanishing gradients can impede training, modern initialization and batch normalization (Ioffe \& Szegedy, 2015) had mitigated the issue for modest depths. Yet the degradation persisted. He et al. (2016) provided convincing evidence with a simple experiment on CIFAR-10: a 56-layer plain network incurred higher training and test errors than a 20-layer counterpart. This pointed to a fundamental optimization difficulty: the stacked nonlinear layers could not easily approximate identity mappings. The key insight of ResNet was to explicitly reformulate the layers to learn a residual function $F(x) = H(x) - x$ instead of directly learning $H(x)$. This seemingly minor change, implemented via shortcut connections that perform identity mapping and element-wise addition, completely altered the optimization landscape.

The contributions of the ResNet paper are manifold. First, it introduced a practical and parameter-free mechanism — residual learning — that allowed networks to reach unprecedented depths (up to 152 layers) while maintaining trainability. Second, the architecture achieved state-of-the-art results on the ILSVRC 2015 classification task with a top-5 error of 3.57\% (single model) and 3.03\% (ensemble). Third, it dominated the competition in object detection and segmentation tasks when used as a backbone. Fourth, the residual learning principle inspired a family of architectures that continue to evolve. This paper provides a deep dive into every aspect of ResNet, from the problem formulation to architectural details and empirical evidence, culminating in a discussion of its lasting legacy.

The remainder of this paper is organized as follows. Section~\ref{sec:litreview} reviews related work on deep network difficulties and existing mitigation strategies. Section~\ref{sec:method} details the residual learning framework and network architecture designs. Section~\ref{sec:experiments} presents experimental results on ImageNet and CIFAR-10, along with ablation studies. Section~\ref{sec:discussion} offers deeper insight into why residual networks work, covering gradient flow, ensemble interpretation, and forward information propagation. Finally, Section~\ref{sec:conclusion} concludes with a summary and discussion of future directions.

\section{Literature Review}
\label{sec:litreview}

\subsection{Early Deep Networks and the Depth Dilemma}
The journey toward deeper networks began with AlexNet (Krizhevsky et al., 2012), which contained 8 layers. VGG (Simonyan \& Zisserman, 2015) demonstrated that increasing depth to 16–19 layers with small 3×3 convolutions consistently improved accuracy. GoogLeNet (Szegedy et al., 2015) employed inception modules and achieved 22 layers while controlling computational cost. However, attempts to go further encountered obstacles. Even with careful initialization and batch normalization, plain networks with 30+ layers showed performance saturation and eventual degradation. This was not simply due to overfitting — the training error itself increased, suggesting that the optimization algorithm struggled to find good minima. Researchers explored various regularizations and auxiliary classifiers, but the degradation remained a bottleneck.

\subsection{Vanishing/Exploding Gradients and Normalization}
Earlier recognition that very deep networks suffer from vanishing or exploding gradients led to techniques such as layer-wise pre-training (Hinton \& Salakhutdinov, 2006) and gradient clipping. The advent of batch normalization (Ioffe \& Szegedy, 2015) significantly improved gradient flow by normalizing layer inputs, enabling training of up to 30-layer plain networks. Yet even with batch normalization, deeper stacks exhibited degradation. This confirmed that the issue was not merely gradient magnitude but the difficulty of learning identity mappings across many nonlinear transformations.

\subsection{Prior Attempts to Improve Depth Training}
Several works aimed to ease depth-related difficulties. Highway networks (Srivastava et al., 2015) introduced gated shortcut connections that allowed information to flow across layers. However, they required learnable gates, adding parameters and complexity. In contrast, ResNet's shortcut connections are simple, parameter-free identity mappings, making them both elegant and effective. Other approaches like deeply-supervised nets (Lee et al., 2015) inserted auxiliary classifiers to provide additional gradient signals, but this increased training overhead. The residual learning framework offered a clean solution that could be integrated into any existing convolutional architecture.

\subsection{Impact of the ResNet Paper}
The ResNet paper (He et al., 2016) not only solved the degradation problem but also set new benchmarks across multiple visual tasks. Its impact extended beyond classification: object detection (Faster R-CNN with ResNet), semantic segmentation (DeepLab), and even video analysis adopted ResNet backbones. The paper's ablation studies — comparing identity shortcuts versus projection shortcuts versus zero-padding — provided rigorous justification for design choices. Subsequent work explored pre-activation variants (He et al., 2016b), wider residual blocks (Zagoruyko \& Komodakis, 2016), and aggregated transformations (Xie et al., 2017), all tracing their lineage to the original ResNet.

\section{Methodology}
\label{sec:method}

\subsection{Residual Learning Formulation}
Consider a stack of layers that should, ideally, map an input $x$ to $H(x)$. If the optimal mapping is close to the identity transformation (as tends to happen when depth is excessive), it is easier to learn the deviation $F(x) = H(x) - x$ than the entire $H(x)$. In residual learning, the layer outputs are added to their input via a shortcut connection:
\begin{equation}
    y = \mathcal{F}(x, \{W_i\}) + x,
    \label{eq:residual}
\end{equation}
where $y$ is the output, $\mathcal{F}(x, \{W_i\})$ represents the residual mapping to be learned (typically two or three convolutional layers with batch normalization and ReLU), and $x$ is the input. The addition is element-wise. If the dimensions of $x$ and $\mathcal{F}$ differ (e.g., when changing the number of channels or reducing spatial size), a linear projection $W_s$ can be used to match dimensions:
\begin{equation}
    y = \mathcal{F}(x, \{W_i\}) + W_s x.
    \label{eq:proj}
\end{equation}
However, the original paper found that identity shortcuts ($W_s = I$) are sufficient for most cases and yield the best results; projection shortcuts only add parameters without significant improvement.

\subsection{Why Residuals Facilitate Optimization}
The key insight is that when the optimal mapping $H^*(x)$ is close to identity, the residual function $\mathcal{F}^*(x) = H^*(x) - x$ has small magnitude and is easier to learn. Mathematically, the gradient of the loss $\mathcal{L}$ with respect to the input of a residual block can be written as:
\begin{equation}
    \frac{\partial \mathcal{L}}{\partial x} = \frac{\partial \mathcal{L}}{\partial y} \left(1 + \frac{\partial \mathcal{F}}{\partial x}\right).
    \label{eq:gradient}
\end{equation}
The presence of the term $1$ ensures that the gradient can flow directly to earlier layers without being scaled by repeated multiplications through nonlinearities. This prevents gradient vanishing even in very deep networks. Furthermore, the residual mapping $\mathcal{F}$ only needs to learn perturbations, which are typically small and easier to optimize.

\subsection{Shortcut Connection Variants}
He et al. (2016) experimentally compared three types of shortcuts:
\begin{itemize}
    \item \textbf{A: Identity shortcut with zero-padding} — when dimensions increase, use zero-padded extra channels (no extra parameters).
    \item \textbf{B: Projection shortcut} — use 1×1 convolutions to match dimensions, applied only when sizes change.
    \item \textbf{C: All shortcuts are projections} — every block uses 1×1 convolutions for shortcuts.
\end{itemize}
On ImageNet, option A performed comparably to B but without additional parameters. Option C was slightly better than B but added computational cost. The identity shortcut (option A) became the standard due to its simplicity and effectiveness.

\subsection{Network Architecture}
\label{subsec:architecture}

ResNet builds upon the VGG spirit of using only 3×3 convolutions but introduces residual blocks. Two basic building blocks are defined:
\begin{itemize}
    \item \textbf{Basic Block}: two consecutive 3×3 convolutions, each followed by batch normalization and ReLU (except the last ReLU after addition). Used in ResNet-18 and ResNet-34.
    \item \textbf{Bottleneck Block}: three layers — a 1×1 convolution to reduce channels, a 3×3 convolution, and another 1×1 convolution to restore channels. This design reduces parameter count while maintaining representational power. Used in deeper variants (ResNet-50/101/152).
\end{itemize}

Table~\ref{tab:arch} lists the architecture configurations for widely used ResNet variants. Note that the spatial resolution is downsampled by convolutional layers with stride 2 at several stages (conv3\_1, conv4\_1, conv5\_1).

\begin{table}[ht]
\centering
\caption{ResNet architectures for ImageNet. The number of output channels is shown for each stage. The final average pooling and 1000-way fully connected layer are common.}
\label{tab:arch}
\begin{tabular}{c|c|c|c|c}
\toprule
Layer Name & Output Size & ResNet-18 & ResNet-34 & ResNet-50 \\
\midrule
conv1 & 112×112 & 7×7, 64, stride 2 & same & same \\
conv2\_x & 56×56 & $\begin{bmatrix} 3\times3, 64 \\ 3\times3, 64 \end{bmatrix} \times 2$ & $\begin{bmatrix} 3\times3, 64 \\ 3\times3, 64 \end{bmatrix} \times 3$ & $\begin{bmatrix} 1\times1, 64 \\ 3\times3, 64 \\ 1\times1, 256 \end{bmatrix} \times 3$ \\
conv3\_x & 28×28 & $\begin{bmatrix} 3\times3, 128 \\ 3\times3, 128 \end{bmatrix} \times 2$ & $\begin{bmatrix} 3\times3, 128 \\ 3\times3, 128 \end{bmatrix} \times 4$ & $\begin{bmatrix} 1\times1, 128 \\ 3\times3, 128 \\ 1\times1, 512 \end{bmatrix} \times 4$ \\
conv4\_x & 14×14 & $\begin{bmatrix} 3\times3, 256 \\ 3\times3, 256 \end{bmatrix} \times 2$ & $\begin{bmatrix} 3\times3, 256 \\ 3\times3, 256 \end{bmatrix} \times 6$ & $\begin{bmatrix} 1\times1, 256 \\ 3\times3, 256 \\ 1\times1, 1024 \end{bmatrix} \times 6$ \\
conv5\_x & 7×7 & $\begin{bmatrix} 3\times3, 512 \\ 3\times3, 512 \end{bmatrix} \times 2$ & $\begin{bmatrix} 3\times3, 512 \\ 3\times3, 512 \end{bmatrix} \times 3$ & $\begin{bmatrix} 1\times1, 512 \\ 3\times3, 512 \\ 1\times1, 2048 \end{bmatrix} \times 3$ \\
\midrule
FLOPs & & $1.8\times10^9$ & $3.6\times10^9$ & $3.8\times10^9$ \\
\bottomrule
\end{tabular}
\end{table}

\subsection{Implementation Details for Training}
All ResNet models are trained from scratch using stochastic gradient descent (SGD) with momentum 0.9, weight decay 1e-4, and mini-batch size 256. Learning rate starts at 0.1 and is divided by 10 when error plateaus. For ImageNet, standard data augmentation (random crop, horizontal flip, color jitter) is used. Batch normalization is applied right after each convolution and before ReLU, following the pre-activation practice that improved stability (He et al., 2016b).

\section{Experiments and Results}
\label{sec:experiments}

\subsection{ImageNet Classification}
The main experiments were conducted on the ILSVRC 2012 dataset with 1.28 million training images and 50,000 validation images. He et al. compared plain and residual versions of 18-layer and 34-layer networks. The 34-layer plain network exhibited a clear degradation: its training error was consistently higher than that of the 18-layer plain network throughout training. In contrast, the 34-layer ResNet not only matched but outperformed the 18-layer ResNet, showing no degradation. Table~\ref{tab:imagenet_error} summarizes the top-1 and top-5 validation errors.

\begin{table}[ht]
\centering
\caption{Single-model results on ImageNet validation set (center-crop evaluation).}
\label{tab:imagenet_error}
\begin{tabular}{lccc}
\toprule
Method & Top-1 error (\%) & Top-5 error (\%) & Depth \\
\midrule
VGG-16 & 28.07 & 9.33 & 16 \\
GoogLeNet & - & 9.15 & 22 \\
Plain-34 & 28.54 & 12.18 & 34 \\
ResNet-34 & 25.03 & 7.76 & 34 \\
ResNet-50 & 22.85 & 6.28 & 50 \\
ResNet-101 & 21.75 & 5.71 & 101 \\
ResNet-152 & 21.43 & 5.71 (3.57 ensemble) & 152 \\
\bottomrule
\end{tabular}
\end{table}

The 152-layer ResNet achieved a top-5 error of 3.57\% (single model), significantly surpassing the previous best of 6.67\% from GoogLeNet. When combining multiple models, the ensemble reached 3.03\%, winning the ILSVRC 2015 classification task. ResNet also demonstrated efficiency: despite being 8 times deeper than VGG-19, it had fewer parameters (25.5M for ResNet-152 vs 144M for VGG-19) and lower computational cost.

\subsection{CIFAR-10 Experiments}
To further investigate depths beyond ImageNet practical limits, He et al. trained ResNets with up to 1202 layers on CIFAR-10 (50,000 training images, 10 classes). The input images were 32×32, and the network used a stem of 3×3 convolutions followed by three stages of residual blocks. Table~\ref{tab:cifar} shows the results. The 110-layer ResNet achieved a test error of 6.61\%, which was competitive with state-of-the-art at that time. The 1202-layer network still performed reasonably (7.93\%), indicating that residual learning allowed training of extremely deep networks, though with minor overfitting due to the large model relative to dataset size.

\begin{table}[ht]
\centering
\caption{Test errors on CIFAR-10 for various ResNet depths.}
\label{tab:cifar}
\begin{tabular}{lc}
\toprule
Depth & Test error (\%) \\
\midrule
20 & 8.75 \\
32 & 7.51 \\
56 & 6.97 \\
110 & 6.61 \\
1202 & 7.93 \\
\bottomrule
\end{tabular}
\end{table}

\subsection{Ablation and Analysis}
\subsubsection{Shortcut Types}
As mentioned in Section~\ref{sec:method}, three shortcut options were compared on ImageNet with ResNet-34. Option A (identity shortcut with zero-padding) gave top-1 error 25.03\%; Option B (projection shortcut only when dimensions change) gave 24.88\%; Option C (all projections) gave 24.84\%. The differences were small, validating the identity shortcut as a simple and effective choice.

\subsubsection{Bottleneck vs Basic Blocks}
He et al. showed that bottleneck blocks allow deeper networks without excessive parameter growth. ResNet-50 with bottleneck blocks has 25.5M parameters, while a hypothetical residual version using basic blocks for 50 layers would have significantly more parameters. The trade-off between depth and width is thus managed efficiently.

\subsubsection{Pre-Activation Variant}
In a follow-up (He et al., 2016b), the order of operations was changed to batch normalization and ReLU before the convolution (pre-activation). This further improved optimization and allowed training of 1000-layer networks, though performance gains on standard benchmarks were marginal.

\section{Discussion}
\label{sec:discussion}

\subsection{Gradient Flow and Optimization}
Equation~\ref{eq:gradient} showed that the gradient from the loss passes through the shortcut term $1$, allowing error signals to propagate directly to earlier layers. This is in stark contrast to plain networks where every layer multiplies the gradient by the Jacobian of the nonlinearities. Even with batch normalization, plain networks of 30+ layers suffer from gradient degradation. Empirical evidence from gradient histograms confirms that ResNet maintains healthy gradient magnitudes throughout the network, while plain networks show diminishing gradients in early layers.

\subsection{Ensemble Interpretation}
Another viewpoint characterizes ResNets as an implicit ensemble of shallow networks. Because the shortcut connections allow information to skip layers, the effective path length for any particular input can vary. Veit et al. (2016) showed that removing individual residual blocks from a trained ResNet does not drastically hurt performance; actually, most paths are relatively short. This suggests that ResNets behave like a collection of networks of varying depths, each contributing to the final prediction. The unrolled structure of residual blocks creates a multiplicative combination of paths, providing robustness and facilitating training.

\subsection{Forward Information Propagation}
Residual connections also aid forward information flow: identity shortcuts allow low-level features to be preserved and reused in deeper layers without distortion. In image recognition, retaining spatial details across many convolutional layers is crucial; the shortcut connections ensure that high-resolution information can flow through the network without being destroyed by successive bottlenecks. This aligns with the observation that ResNets often produce sharper feature maps in intermediate layers compared to plain networks.

\subsection{Limitations and Improvements}
While ResNet solved the degradation problem, it is not a silver bullet. Residual networks can still be difficult to train if the shortcut path is blocked (e.g., by aggressive downsampling). The pre-activation variant improved gradient propagation by moving the activation outside the residual path. Additionally, batch normalization requires careful tuning of per-layer parameters. Later architectures like DenseNet (Huang et al., 2017) and ResNeXt (Xie et al., 2017) further improved parameter efficiency and performance by modifying the connection patterns.

\section{Subsequent Impact and Extensions}
\label{sec:extensions}

ResNet's influence on deep learning is profound. The residual learning paradigm has become a standard component in nearly all modern architectures. Key subsequent developments include:
\begin{itemize}
    \item \textbf{ResNeXt}: Introduces grouped convolutions to increase width without exploding parameters (Xie et al., 2017).
    \item \textbf{DenseNet}: Connects each layer to every other layer in a feed-forward fashion, encouraging feature reuse and reducing vanishing gradients (Huang et al., 2017).
    \item \textbf{Wide ResNet}: Makes residual blocks wider (more channels) and shallower (fewer layers), achieving better performance with fewer layers (Zagoruyko \& Komodakis, 2016).
    \item \textbf{Pre-activation ResNet}: Changes the order of batch normalization and ReLU to before convolution, enabling 1000-layer training (He et al., 2016b).
    \item \textbf{Attention mechanisms}: Modules like SENet (Hu et al., 2018) and CBAM (Woo et al., 2018) introduced channel and spatial attention on top of residual blocks.
\end{itemize}

Beyond computer vision, ResNet has been adopted in natural language processing (e.g., deep Transformer encoders borrow residual connections), speech recognition (deep CNNs for raw audio), reinforcement learning (deep Q-networks with residual blocks), and generative models. Its legacy lies in the simple yet transformative idea that making the learning target a residual rather than the full mapping dramatically improves optimization.

\section{Conclusion}
\label{sec:conclusion}

The residual learning framework presented by He et al. (2016) directly addressed the degradation problem that plagued deep plain networks. By reformulating block outputs as $F(x) + x$, where $F$ is a residual function, ResNet enabled the training of networks with unprecedented depth (152 layers on ImageNet, 1202 layers on CIFAR-10) while achieving state-of-the-art accuracy. The shortcut connections require no additional parameters, are computationally cheap, and integrate seamlessly with standard CNN layers. Through rigorous experiments and theoretical analysis, the paper demonstrated that residual learning eases optimization by allowing gradients to flow directly through identity paths and by making it easier to learn identity-like mappings.

ResNet's contributions extend far beyond the 2015 ILSVRC results. It established a design principle that has become ubiquitous in deep learning. The architecture's success inspired a wealth of variants and cross-domain applications, cementing its status as one of the most influential papers in the field. The central lesson — that depth is beneficial, but residual connections make deep networks trainable — continues to guide modern research. As the field moves toward even larger models and new architectures, the residual learning concept remains a foundational tool.

\begin{thebibliography}{99}

\bibitem{he2016resnet}
K. He, X. Zhang, S. Ren, and J. Sun.
\newblock Deep residual learning for image recognition.
\newblock In \emph{Proceedings of the IEEE Conference on Computer Vision and Pattern Recognition (CVPR)}, 2016.

\bibitem{simonyan2014vgg}
K. Simonyan and A. Zisserman.
\newblock Very deep convolutional networks for large-scale image recognition.
\newblock In \emph{International Conference on Learning Representations (ICLR)}, 2015.

\bibitem{szegedy2015googlenet}
C. Szegedy, W. Liu, Y. Jia, P. Sermanet, S. Reed, D. Anguelov, D. Erhan, V. Vanhoucke, and A. Rabinovich.
\newblock Going deeper with convolutions.
\newblock In \emph{Proceedings of the IEEE Conference on Computer Vision and Pattern Recognition (CVPR)}, 2015.

\bibitem{krizhevsky2012alexnet}
A. Krizhevsky, I. Sutskever, and G. E. Hinton.
\newblock Imagenet classification with deep convolutional neural networks.
\newblock In \emph{Advances in Neural Information Processing Systems (NeurIPS)}, 2012.

\bibitem{ioffe2015batchnorm}
S. Ioffe and C. Szegedy.
\newblock Batch normalization: Accelerating deep network training by reducing internal covariate shift.
\newblock In \emph{International Conference on Machine Learning (ICML)}, 2015.

\bibitem{srivastava2015highway}
R. K. Srivastava, K. Greff, and J. Schmidhuber.
\newblock Highway networks.
\newblock arXiv preprint arXiv:1505.00387, 2015.

\bibitem{he2016preact}
K. He, X. Zhang, S. Ren, and J. Sun.
\newblock Identity mappings in deep residual networks.
\newblock In \emph{European Conference on Computer Vision (ECCV)}, 2016.

\bibitem{veit2016ensemble}
A. Veit, M. Wilber, and S. Belongie.
\newblock Residual networks behave like ensembles of relatively shallow networks.
\newblock In \emph{Advances in Neural Information Processing Systems (NeurIPS)}, 2016.

\bibitem{xie2017resnext}
S. Xie, R. Girshick, P. Dollár, Z. Tu, and K. He.
\newblock Aggregated residual transformations for deep neural networks.
\newblock In \emph{Proceedings of the IEEE Conference on Computer Vision and Pattern Recognition (CVPR)}, 2017.

\bibitem{huang2017densenet}
G. Huang, Z. Liu, L. van der Maaten, and K. Q. Weinberger.
\newblock Densely connected convolutional networks.
\newblock In \emph{Proceedings of the IEEE Conference on Computer Vision and Pattern Recognition (CVPR)}, 2017.

\bibitem{zagoruyko2016wideresnet}
S. Zagoruyko and N. Komodakis.
\newblock Wide residual networks.
\newblock In \emph{British Machine Vision Conference (BMVC)}, 2016.

\bibitem{hu2018senet}
J. Hu, L. Shen, and G. Sun.
\newblock Squeeze-and-excitation networks.
\newblock In \emph{Proceedings of the IEEE Conference on Computer Vision and Pattern Recognition (CVPR)}, 2018.

\end{thebibliography}

\appendix

\section{Figure Placeholders}
\label{app:figures}

\begin{figure}[ht]
\centering
\fbox{\parbox{0.8\textwidth}{\centering \vspace{3cm} \textbf{Figure 1:} Training error curves for plain and residual networks on CIFAR-10. The 56-layer plain network exhibits higher error than the 20-layer plain network, illustrating the degradation problem. Residual networks do not suffer from this issue.}}
\caption{Illustration of the degradation problem.}
\label{fig:degradation}
\end{figure}

\begin{figure}[ht]
\centering
\fbox{\parbox{0.8\textwidth}{\centering \vspace{3cm} \textbf{Figure 2:} Architecture of a residual block with identity shortcut (left) and a bottleneck block (right). The identity shortcut adds the input to the output; the bottleneck uses 1×1 convolutions to reduce and expand channels.}}
\caption{Residual block structures.}
\label{fig:blocks}
\end{figure}

\end{document}
